\subsection{How to Read Matrix Problems}

\begin{corebox}[Questions to ask before calculating]
When working with matrices, determine:
\begin{itemize}
    \item the size of every matrix;
    \item which operations are defined;
    \item whether the operation is entrywise or row-by-column;
    \item whether the order of multiplication matters;
    \item whether visible structure simplifies the problem;
    \item what the result says about the original objects.
\end{itemize}
\end{corebox}

\begin{interpretationbox}[A matrix solution is more than its last line]
A complete solution should expose the compatibility checks, the chosen operation,
the main calculation, and the meaning of the result. The purpose is not merely to
produce an array of numbers, but to use the matrix as a mathematical description.
\end{interpretationbox}

\begin{summarybox}[The matrix workflow]
Read the shape, identify the operation, calculate according to its definition,
and interpret the resulting structure. This workflow prepares the later use of
matrices in determinants, inverses, and systems of equations.
\end{summarybox}
